\documentclass[11pt]{article}
\usepackage[a4paper,margin=1in]{geometry}
\usepackage{hyperref}
\usepackage{longtable}
\usepackage{array}
\usepackage{enumitem}
\usepackage{xcolor}
\usepackage{graphicx}
\usepackage{float}
\usepackage{booktabs}

\hypersetup{
  colorlinks=true,
  linkcolor=blue!60!black,
  urlcolor=blue!60!black,
  citecolor=blue!60!black
}

\setlist[itemize]{leftmargin=1.2em}
\setlist[enumerate]{leftmargin=1.5em}

\begin{document}

\begin{titlepage}
  \centering
  {\LARGE Req2Design \par}
  \vspace{0.5cm}
  {\Large Software Requirements Specification (SRS) Report\par}
  \vspace{1cm}
  {\large Author: ali \par}
  \vfill
  {\large \today\par}
  \vspace{2cm}
  \textbf{Overview}\par
  This report captures the Req2Design SRS, summarizes the current implementation (React + Flask requirements-engineering platform), traces requirements to the existing codebase, and outlines a test plan with test cases and current execution status for Overleaf sharing.
\end{titlepage}

\tableofcontents
\newpage

\section{Introduction}
\subsection{Purpose}
Define the requirements for Req2Design, a web-based system that ingests vague requirements (text or audio), validates/sanitizes them, and generates IEEE-830-style SRS artifacts, textual use cases, and use-case diagrams.

\subsection{Scope}
The SRS covers functional and non-functional requirements for the requirements-to-SRS platform. Detailed design is summarized only to show alignment between the platform and the SRS needs.

\subsection{Definitions and Acronyms}
\begin{itemize}
  \item UI: User Interface
  \item API: Application Programming Interface
  \item DBMS: Database Management System
  \item SSL: Secure Sockets Layer
\end{itemize}

\subsection{References}
\begin{itemize}
  \item IEEE Std 830-1998: Software Requirements Specification
  \item Project README and backend source code in \texttt{api\_server.py}, \texttt{main\_orchestrator.py}, and \texttt{srs\_model\_generator.py}
\end{itemize}

\subsection{Overview}
Req2Design is a browser-based application for capturing raw requirements and producing structured SRS outputs with export to JSON/HTML/PDF.

\section{Overall Description}
\subsection{Product Perspective}
Target platform: responsive web app with a Flask backend and React frontend. The platform ingests requirements (text/audio), validates them, and uses LLM-backed generation to produce SRS outputs.

\subsection{Product Functions}
\begin{itemize}
  \item Capture requirements via text input and audio upload.
  \item Validate and sanitize inputs to mitigate prompt-injection or spam.
  \item Generate IEEE 830-style SRS JSON, textual use cases, and diagram data.
  \item Export SRS to JSON/HTML/PDF for sharing.
  \item Provide light/dark UI theme toggle.
  \item Support authentication-ready flows (frontend scaffolding in place).
  \item Maintain processing history and optional persistence hooks.
\end{itemize}

\subsection{User Characteristics}
General public, ages 18--65, basic-to-intermediate technical proficiency, English language.

\subsection{Constraints}
\begin{itemize}
  \item Work across major browsers (Chrome, Firefox, Safari, Edge).
  \item Secure API communication over HTTPS.
  \item External dependency on LLM provider (Replicate/HF) for generation.
\end{itemize}

\subsection{Assumptions and Dependencies}
Assumes availability of an LLM API token (HF/Replicate), optional RDBMS for persistence, and HTTPS termination.

\section{Specific Requirements}
\subsection{External Interfaces}
\begin{itemize}
  \item \textbf{UI:} Responsive web experience (desktop/mobile).
  \item \textbf{Hardware:} Standard desktop/laptop.
  \item \textbf{Software:} Chrome, Firefox, Safari, Edge on Windows/macOS/Linux.
  \item \textbf{Communication:} HTTPS; LLM provider APIs for generation.
\end{itemize}

\subsection{Functional Requirements}
\begin{longtable}{p{0.13\linewidth} p{0.62\linewidth} p{0.15\linewidth}}
\toprule
\textbf{ID} & \textbf{Description} & \textbf{Priority} \\
\midrule
FR-1 & Accept text-based requirements submission & High \\
FR-2 & Accept audio upload and transcribe to text & High \\
FR-3 & Sanitize and validate inputs (length, injection, repetition) & High \\
FR-4 & Generate SRS JSON (IEEE 830 style) from inputs & High \\
FR-5 & Generate textual use cases and use-case diagram data & High \\
FR-6 & Export SRS to JSON/HTML/PDF & High \\
FR-7 & Provide light/dark theme toggle in UI & Medium \\
FR-8 & Allow basic user auth scaffolding (login/signup views) & Medium \\
FR-9 & Maintain processing history and optional DB persistence & Medium \\
FR-10 & Provide health/status endpoints for monitoring & Medium \\
\bottomrule
\end{longtable}

\subsection{Non-Functional Requirements}
\begin{itemize}
  \item \textbf{Performance:} Typical request/response $\leq$ 5 seconds with LLM latency budgeted; handle bursty loads when queued.
  \item \textbf{Security:} TLS for data in transit; prompt-injection defenses on inputs; secure handling of API tokens.
  \item \textbf{Reliability/Availability:} Target $<$1\% monthly downtime; graceful degradation if LLM unavailable.
  \item \textbf{Usability:} Simple UI, clear feedback, accessible theme toggle.
  \item \textbf{Maintainability:} Modular backend (Flask + orchestrator) and frontend components; logging and configuration separation.
  \item \textbf{Portability:} Deployable on common OS/cloud with Python and Node runtimes.
\end{itemize}

\section{Implementation Overview (Current Repository)}
\subsection{Architecture}
\begin{itemize}
  \item \textbf{Frontend:} React + Tailwind (\texttt{frontend/}) for requirements input, history, and SRS viewing.
  \item \textbf{Backend:} Flask API in \texttt{api\_server.py} with endpoints for health, single/batch processing, audio transcription, SRS generation, download, stats, and cleanup.
  \item \textbf{Processing Core:} \texttt{main\_orchestrator.py} wires the requirements processor, batch processor, SRS exporters, and model-driven SRS generator.
  \item \textbf{Model Layer:} \texttt{srs\_model\_generator.py} uses Replicate-hosted LLM to generate SRS content; \texttt{srs\_generator.py} exports JSON/HTML/PDF.
  \item \textbf{Persistence:} Optional database hooks via \texttt{requirements.db} and stubs for \texttt{DataManager}; file outputs in \texttt{data/output}.
\end{itemize}

\subsection{Key API Endpoints}
\begin{itemize}
  \item \texttt{GET /api/health} --- service heartbeat.
  \item \texttt{POST /api/process-single} --- process text requirement.
  \item \texttt{POST /api/process-audio} --- process uploaded audio.
  \item \texttt{POST /api/transcribe-audio} --- transcribe audio only.
  \item \texttt{POST /api/process-batch} --- batch file processing.
  \item \texttt{POST /api/generate-srs} --- generate SRS from processed results.
  \item \texttt{POST /api/download-srs/\{format\}} --- export SRS in JSON/HTML/PDF.
  \item \texttt{POST /api/generate-srs-pdf} --- convenience SRS PDF.
  \item \texttt{POST /api/generate-srs-from-audio} --- audio-to-SRS shortcut.
  \item \texttt{POST /api/generate-srs-from-file} --- file-to-SRS shortcut.
  \item \texttt{GET /api/stats} --- system stats.
  \item \texttt{POST /api/cleanup} --- remove temp files.
\end{itemize}

\subsection{Input Handling and Security}
\begin{itemize}
  \item Prompt-injection detection via regex filters and sanitization for text and audio transcripts.
  \item Validation: minimum word count, alphabetic content check, repetition detection, and length truncation.
  \item Logging: warnings for suspicious input and sanitization deltas.
\end{itemize}

\section{Requirements Traceability (SRS $\rightarrow$ Implementation)}
The current codebase already fulfills most platform-oriented requirements for requirements-to-SRS generation. The table maps SRS needs to existing modules to guide coverage and future improvements.

\begin{longtable}{p{0.13\linewidth} p{0.42\linewidth} p{0.28\linewidth} p{0.12\linewidth}}
\toprule
\textbf{ID} & \textbf{Required Behavior} & \textbf{Current/Planned Module} & \textbf{Status} \\
\midrule
FR-1 & Text submission & React frontend forms + \texttt{/api/process-single} & Implemented \\
FR-2 & Audio submission/transcription & React upload + \texttt{/api/process-audio} & Implemented \\
FR-3 & Input sanitization/validation & \texttt{sanitize\_user\_input}, \texttt{validate\_text\_content} & Implemented \\
FR-4 & Generate SRS JSON & \texttt{srs\_model\_generator.py} + \texttt{/api/generate-srs} & Implemented (LLM-dependent) \\
FR-5 & Use cases / diagram data & Model output path in \texttt{srs\_model\_generator.py} & Implemented (LLM-dependent) \\
FR-6 & Export JSON/HTML/PDF & \texttt{srs\_generator.py}, \texttt{/api/download-srs/\{format\}} & Implemented \\
FR-7 & Theme toggle & Header component (light/dark) & Implemented \\
FR-8 & Auth scaffolding & Frontend login/signup components & Implemented (stub) \\
FR-9 & Processing history/persistence & \texttt{DataManager} stubs + \texttt{requirements.db} hook & Partial (stub) \\
FR-10 & Health/status endpoints & \texttt{/api/health}, \texttt{/api/stats} & Implemented \\
\bottomrule
\end{longtable}

\section{Test Plan, Cases, and Current Results}
\subsection{Strategy}
Combine automated API tests with manual UI checks. Security tests cover input sanitization and prompt-injection detection. Performance tests target latency and concurrency once travel services exist.

\subsection{Test Environment}
\begin{itemize}
  \item Backend: Python 3.10+, Flask, Replicate token configured.
  \item Frontend: Node 18+, React dev server.
  \item Data: sample requirements in \texttt{data/input}; generated SRS in \texttt{data/output}.
\end{itemize}

\subsection{Test Cases}
\begin{longtable}{p{0.16\linewidth} p{0.45\linewidth} p{0.18\linewidth} p{0.14\linewidth}}
\toprule
\textbf{ID} & \textbf{Scenario} & \textbf{Expected Result} & \textbf{Status} \\
\midrule
TC-API-1 & \texttt{/api/health} returns service heartbeat & JSON with \texttt{status=healthy} & Pass \\
TC-API-2 & Process valid text requirement via \texttt{/api/process-single} & Returns processed requirement payload & Not run (token required) \\
TC-API-3 & Reject prompt-injection text & 4xx with validation error & Not run \\
TC-API-4 & Transcribe audio via \texttt{/api/process-audio} & JSON with transcript and processing result & Not run (token + audio sample) \\
TC-API-5 & Generate SRS from processed batch & SRS JSON/HTML/PDF saved to output dir & Not run \\
TC-SEC-1 & Detect repeated-token spam & Validation fails with error & Not run \\
TC-PERF-1 & Process 100 concurrent requests & p95 latency $\leq$ 5s & Not run \\
TC-UX-1 & Frontend requirements submission flow & User can submit text and view SRS & Not run \\
\bottomrule
\end{longtable}

\subsection{How to Execute}
\begin{enumerate}
  \item Configure \texttt{HF\_API\_TOKEN} or \texttt{REPLICATE\_API\_TOKEN} as required.
  \item Start backend: \texttt{python api\_server.py}.
  \item Start frontend: \texttt{cd frontend \&\& npm start}.
  \item Use cURL/Postman for API cases; capture outputs in \texttt{data/output} and paste summaries back into this table.
\end{enumerate}

\section{Documentation and Overleaf Usage}
\begin{itemize}
  \item Place this file in Overleaf as \texttt{main.tex} or upload the repository with a \texttt{docs/} folder.
  \item Add any generated SRS JSON/HTML/PDF from \texttt{data/output} as appendices or artifacts.
  \item Keep code documentation in-line (docstrings/comments) and mirror key design notes here for reviewers.
\end{itemize}

\section{Future Work to Meet SRS Fully}
\begin{itemize}
  \item Harden persistence: replace stubs with full DB layer and migrations.
  \item Expand automated tests for all endpoints and UI flows; add CI.
  \item Add role-based authN/Z, rate limiting, and monitoring for production readiness.
  \item Improve resilience to upstream LLM failures with retries/fallbacks.
\end{itemize}

\end{document}

